%% kapa: exam=1 round=30 year=2019 period=2 form=A subject=appraisal_law subject_ko=감정평가관계법규 seq=33 number=33
%% source: Q-Net 감정평가사 기출문제 (cst00309) artlSeq=5205205 file=extracted/1008/5205205/2019년 제30회 감정평가사 1차시험 문제지/2019년 제30회 감정평가사 1차시험 문제지/2019년도 제30회 감정평가사 1차 2교시 A형.hwp
%% conversion: kordoc 4.24.51 -> tools/gen1.py (auto); review=proofread-sonnet; flags=none
\begin{question}{33}
부동산등기법령상 수용으로 인한 등기에 관한 설명으로 옳은 것은?
\begin{choices}
  \item 수용으로 인한 소유권이전등기는 등기권리자가 단독으로 신청할 수 없다.
  \item 등기관이 수용으로 인한 소유권이전등기를 하는 경우 그 부동산의 등기기록 중 소유권 외의 권리에 관한 등기가 있으면 그 등기를 당사자의 신청에 따라 말소하여야 한다.
  \item 부동산에 관한 소유권 외의 권리의 수용으로 인한 권리이전등기에 관하여는 수용으로 인한 소유권이전등기 규정이 적용되지 않는다.
  \item 등기관이 수용으로 인한 소유권이전등기를 하는 경우 그 부동산의 등기기록 중 그 부동산을 위하여 존재하는 지역권의 등기는 직권으로 말소할 수 없다.
  \item 수용으로 인한 소유권이전등기를 신청하는 경우에 등기권리자는 포괄승계인을 갈음하여 포괄승계로 인한 소유권이전의 등기를 신청할 수 없다.
\end{choices}
\end{question}
\answer{4}
